\documentclass[10pt,a4paper]{article}
\usepackage{fontspec,xeCJK}
\setCJKmainfont{HaranoAjiMincho-Regular.otf}[BoldFont=HaranoAjiMincho-Bold.otf]
\usepackage{amsmath,amssymb,amsthm,booktabs,geometry,hyperref,xurl}
\geometry{margin=24mm}
\newcommand{\KF}{K_{\mathrm F}}
\newcommand{\conv}{\operatorname{conv}}
\newcommand{\eps}{\varepsilon}
\newtheorem{theorem}{定理}
\newtheorem{proposition}[theorem]{命題}
\newtheorem{lemma}[theorem]{補題}
\renewcommand{\proofname}{証明}
\title{一般化 T 区間と Hall's ray より下の Markov 区間\\\large 31313禁止による区間構成と131禁止の課題}
\author{}
\date{2026年10月1日}
\begin{document}
\maketitle
\begin{abstract}
Freiman の31313禁止の閉じた許容表を、新しい固定語の対 $(322,431)$ に接続し、
$[4.52578,4.52754]\subset M\cap(-\infty,c_F)$ を証明する。
許容表の全3,464,816行を固定したまま後続被覆を再検証し、
全ての非中心値を一様に抑えることで、Cantor 和の区間を Markov 値の区間に移す。
別実装による有理数だけの非中心上界の検算も行った。
一方、当初の131禁止集合 $\KF+\KF$ の内点構成は未完了である。
後半では、全接頭語版の一般化 T 区間の定義、証明済みの障害、探索の到達範囲を区別して記す。
\end{abstract}

\section{Hall's ray より下の Markov 区間：完成した構成}\label{sec:below-ray}
本節では、後の131禁止集合とは区別して、Freiman フォルダーの
\emph{31313禁止}の閉じた許容族を用いる。この変更によって
$M$ の Hall's ray 以外の内点という目標は達成できる。
$\KF+\KF$ の内点を示したという主張ではない。

\subsection{具体的な区間}
有限語 $a$ の固定部分には4を許し、追加する無限語の数字は1,2,3に限る。
\[
 \widetilde K(a)=\{[0;a,w]:w\in\{1,2,3\}^{\mathbb N},\ aw\text{ が31313を避ける}\}.
\]
特に固定語431の4は許される。これは、固定部分にも4の禁止を課す
後述の $C_{\{4,131\}}(a)$ と異なる規約である。
\[
 H_{a,b}=\conv(\widetilde K(a)+\widetilde K(b))=[L_{a,b},H_{a,b}^{+}],\qquad
 J_{p,q}(a,b)=[L_{a,b}+p(H_{a,b}^{+}-L_{a,b}),\,
 L_{a,b}+q(H_{a,b}^{+}-L_{a,b})].
\]
以下では $a=322,b=431$ とし、凸包の端点を単に $L,H$ と書く。
交互辞書式順序による極値計算から
\[
 L=\frac{8512493+28004\sqrt{462}}{17339617},\qquad
 H=\frac{1352829-8081\sqrt{462}}{2233974}.
\]

\begin{theorem}[Hall's ray より下の正長区間]\label{thm:below-ray}
Freiman の固定許容表の全行再検証と、以下の初期条件の厳密検証により
\[
 4+J_{1/16,\,7/8}(322,431)\subset M
\]
が成立する。この区間はおよそ
\[
 [4.525777278415714\ldots,\ 4.527546990114258\ldots]
\]
であり、特に有理端点を持つ区間について
\[
 \boxed{[4.52578,\ 4.52754]\subset M\cap(-\infty,c_F)}
\]
を得る。ここで
$c_F=(2221564096+283748\sqrt{462})/491993569$ は Hall's ray の左端である。
\end{theorem}

\subsection{閉じた後続被覆への接続}
Freiman の一般族は、31313状態、偶奇、末尾2桁、分母比の形状、微分比 $S$、
区間型 $(p,q)$ によって定義される。
$\xi_a$ を $F_a(\xi_a)=\min\widetilde K(a)$ となる尾とし、
\[
 r=\frac{q_{a^-}}{q_a},\quad s=\frac{q_{b^-}}{q_b},\quad
 S=\left(\frac{q_a}{q_b}\right)^2
       \frac{(1+r\xi_a)^2}{(1+s\xi_b)^2},\qquad \kappa=\frac{51}{50}.
\]
今回の固定語について厳密に
\[
 \begin{gathered}
 \sigma(a)=\eps,\quad\sigma(b)=31,\quad |a|\equiv|b|\equiv1\pmod2,\qquad
 q_a=q_b=17,\\
 r=\frac7{17}\in\left[\frac9{22},\frac{14}{33}\right],\qquad
 s=\frac{13}{17}\in\left[\frac{13}{17},\frac{19}{24}\right],\\
 S=\frac{55792801+2467176\sqrt{462}}{159491641}
   \in[\kappa^{-20},\kappa^{-19}]
 \end{gathered}
\]
が成立する。形状の箱は逆向き末尾22と13に $[1/4,4/5]$ を代入した像である。
Freiman の \texttt{graph\_wide} の採用表には、この状態対と比の箱に対し
\[
 (p,q)=\left(\frac j{32},\frac{j+2}{32}\right),\qquad j=2,3,\ldots,26
\]
の25行がすべて含まれる。隣り合う区間は重なり、その和集合が
$J_{1/16,7/8}(322,431)$ になる。

この表に採用された3,464,816行すべてについて、許容的な真の3後続の有限個の和が
親区間を覆うことを再検証した。具体的には、代数的端点を分母 $2^{48}$ の
有理区間で囲った入力を再生成し、保存入力との完全一致を確かめたうえで、
保存した採用表を\emph{削除・修正せずに}検証器へ渡す。
後続の形状の箱への包含、微分比の像が触れる全ての箱での子の許容性、
区間の重なりの両向きの不等式、親の両端の被覆を確認する。
全行が通り、採用表のバイト列も変わらない。

表全体で $\kappa^{-155}\le S\le\kappa^{155}$ であり、
$1/4\le S/(q_a/q_b)^2\le4$ だから左右の分母比は正の上下界を持つ。
従って、被覆からある数を含む子を選び続けると両側の語長が無限に伸びる。
円筒幅が0に向かい、極限も31313を避けることから
\[
 J_{1/16,7/8}(322,431)\subset\widetilde K(322)+\widetilde K(431).
\]
ここまでは Cantor 和の区間包含であり、次に $M$ への包含を確認する。

\subsection{全ての非中心値の上界}
二側列 $\theta$ の局所値を
\[
 \lambda_n(\theta)=\theta_n+[0;\theta_{n-1},\theta_{n-2},\ldots]
                         +[0;\theta_{n+1},\theta_{n+2},\ldots],\qquad
 m(\theta)=\sup_{n\in\mathbb Z}\lambda_n(\theta)
\]
とする。今、$\lambda_0=4+x+y$、$x\in\widetilde K(322)$、
$y\in\widetilde K(431)$ となる列を考える。固定中心部は
\[
 \cdots\,2\,2\,3\,4^*\,4\,3\,1\,\cdots
\]
である。
\begin{lemma}\label{lem:below-noncentral}
この固定中心部を持つ上の全ての列について
\[
 \sup_{n\ne0}\lambda_n\le
 \Theta:=\frac{4\sqrt{462}}{19}+\frac1{3025}
   =4.525422212239698\ldots
   <4+L+\frac{H-L}{16}
\]
が成立する。
\end{lemma}
\begin{proof}
全ての数字が1,2,3で31313を避ける二側列の局所値の最大は
$M_B=4\sqrt{462}/19$ である。合法な5桁の窓を全列挙し、
左右の極値尾を交互辞書式順序で選んで比較すると得られる。
長さ5の窓を固定すれば、長さ5の禁止語が両側の新しい尾を同時に横断しないため、
この最大化は正当である。

固定部内の中心以外の6箇所は同じ極値計算で直接評価する。
その最大上界は位置1の $4.514022576811153\ldots$ である。
固定部の左右それぞれに続く最初の9桁内にある中心は、合法な追加語を列挙して調べる。
数字1,2での局所値は4未満であり、数字3でも中心方向の隣接数字が2以上なら
4.5未満である。それ以外は両方向の連分数を厳密に最大化する。
この検査で全て $\Theta$ 以下になる。列挙した追加接頭語は左右合計57,054個である。

残る中心は、固定部から外向きに10桁以上離れている。
中心方向の最初の9桁を保ち、それ以後を合法な尾に置き換えると、
列全体を1,2,3による31313禁止列にできる。
長さ9の連分数円筒の幅は $1/q_9^2\le1/F_{10}^2=1/3025$ である。
従ってこの置換による局所値の変化は高々 $1/3025$ であり、上界 $\Theta$ を得る。
全ての代数的比較は $\mathbb Q(\sqrt{462})$ で行う。
さらに別実装では、40桁の極値接頭語と終端 $[0,1]$ の有理外側評価だけを使い、
同じ全非中心値に対する上界 $4.525423$ 未満も確認した。
\end{proof}

定理\ref{thm:below-ray}の証明はこれで完了する。
任意の $t\in4+J_{1/16,7/8}(322,431)$ に対し、後続被覆から
$\lambda_0=t$ の二側列を得る。補題\ref{lem:below-noncentral}から
全ての $n\ne0$ で $\lambda_n<t$ なので、$m(\theta)=t$ である。
最後に有理端点区間の包含と上端 $<c_F$ を代数的に比較する。

\subsection{再現方法と131禁止の問題との区別}
リポジトリの根で次を実行する。Python 標準ライブラリと C++17 を使う。
\begin{verbatim}
python3 -B -S Berstein/src/verify_below_ray.py --full
\end{verbatim}
出力は \texttt{Berstein/data/below\_ray\_verified.json} と
\texttt{below\_ray\_certificate.json}、全行再検証の記録は
\texttt{below\_ray\_kernel\_replay.json} に保存する。
Freiman のソースと採用表を読み取り専用で参照し、生成物は全て Berstein 以下に置く。
通常実行では \texttt{--full} を省き、再検証済みデータとソースのハッシュを照合する。
初期語の探索と、固定した初期語の厳密検証は別プログラムになっている。

この構成は Schecker 型の「許容的な真の3後続による閉じた被覆」という方針を実現する。
ただし、区間は全凸包の部分区間 $J_{p,q}$ であり、指定された一般化 T 区間そのものではない。
また31313禁止は131の出現を許す。従って、この結果をそのまま
$\KF+\KF$ や共有された $3+[1.29288,1.292906]$ に移すことはできない。
以下は、このより強い131禁止の問題について得られた障害と未完了の探索である。

\section{集合と一般化 T 区間}
\[
 \KF=\{[0;a_1,a_2,\ldots]:a_i\in\{1,2,3\},\ 131\text{ が現れない}\}.
\]
語は左右とも中心から外向きに読む。$U=(u_1,\ldots,u_h)$ に対し
\[
 F_U(x)=[0;u_1,\ldots,u_h+x],\qquad
 \KF(U)=\{[0;U,w]\in\KF\}
\]
とおく。有限アルファベットを $\{1,2,3,4\}$ とし、有限禁止語集合 $B$ に対して
\[
 C_B(U)=\{[0;U,w]:Uw\text{ の全体に }B\text{ の語が現れない}\}.
\]
固定語の内部、固定語と尾の境界、尾の内部のすべてを判定する。指定された一般化は
\[
 I(U,V;B_1,B_2)=\conv\bigl((C_{B_1}(U)+C_{B_2}(V))
          \cup(C_{B_2}(U)+C_{B_1}(V))\bigr).
\]
一方の和が空ならその項を捨て、両方空なら区間を未定義（コードでは \texttt{None}）とする。
内点構成では、両方の $B_i$ に $4,131$ を含めれば、各項は $\KF+\KF$ の部分集合である。
ただし、その\emph{凸包}まで真の和に含まれるとは限らない。

元の Schecker の T 区間は、固定語の後ろに新しく課す尾条件を用いて
\[
 \conv\bigl((F_U(C(2))+F_V(C(3)))\cup(F_U(C(3))+F_V(C(2)))\bigr)
\]
と書かれる。ここで $C(j)$ は数字 $1,\ldots,j$ を使う連分数集合である。
したがって、全接頭語に遡る上の規約では、$B_1=\{3,4\},B_2=\{4\}$ は
\emph{すべての固定語について}元の T 区間を再現するわけではない。
例えば $U=V=3$ では全接頭語版の二つの項はともに空だが、元の T 区間は非空である。
この二つの規約を黙って交換してはならない。

全接頭語規約では、現在の131状態と短い末尾だけで、追加禁止条件の可否は判定できない。
例えば11221と13221は同じ偶奇・末尾221・131状態を持つが、追加禁止語13を
前者は避け、後者は既に含んでいる。一般化 T の有限状態表には、各追加禁止条件について
「過去に違反したか」の情報も必要である。直接探索のコードは固定語全体を検査しており、
この情報を捨てていない。後述の幾何的な部分区間探索は、追加禁止語型を使わない別の探索である。

\section{被覆から内点を導くために必要な閉性}
3後続とは、追加語 $u,v$ の数字が $1,2,3$ であるという意味で用いる。
後続の個数を三つに制限する意味ではない。

\begin{theorem}[一般化 T 区間による帰納原理]\label{thm:cover}
有限語対と補助条件からなる許容族 $\mathcal A$ があり、各要素 $P$ に
非空な閉区間 $I_P$ が対応するとする。以下を仮定する。
\begin{enumerate}
\item $U_P,V_P$ は131を避け、$I_P\subseteq H(U_P,V_P)$、ただし
 $H(U,V)=\conv(\KF(U)+\KF(V))$。
\item 各 $P$ には有限個の許容的な子 $Q$ があり、
 $U_Q=U_Pu,V_Q=V_Pv$、$|u|+|v|>0$、追加数字は1,2,3であり、
 $I_P\subseteq\bigcup_Q I_Q$。
\item 許容族全体で $0<c\le |\conv\KF(V_P)|/|\conv\KF(U_P)|\le C<\infty$。
\end{enumerate}
このとき各 $I_P$ は $\KF(U_P)+\KF(V_P)$ に含まれる。
特に一つでも正長の許容区間があれば $\KF+\KF$ は内点を持つ。
\end{theorem}
\begin{proof}
$t\in I_P$ を固定し、被覆により $t$ を含む許容子を反復して選ぶ。
総追加長は無限に増大する。連分数の分母は長さとともに発散し、各円筒幅は0に向かう。
片側だけが無限に伸びると、伸びない側の幅は正の定数にとどまり、仮定3に反する。
よって両側とも無限に伸びる。入れ子の基礎円筒の和の凸包 $H(U_n,V_n)$ は
幅が0に向かい、そのすべてに $t$ が入る。極限の二つの連分数を $x,y$ とすると
$x+y=t$ であり、全段階で131を避けるため $x,y\in\KF$。
\end{proof}

補助の $B_1,B_2$ は子で変更してよい。ただし\emph{基礎の131禁止と接頭語の延長}は保つ。
この証明には、補助区間どうしの入れ子性も、追加禁止条件の単調な強化も必要ない。
反対に、「その親で短い被覆が見つかった」だけでは仮定2の\emph{全許容子への閉性}を満たさない。

\section{131禁止の極値を厳密に求める}
禁止語の接頭語でもある最長末尾を状態とする。
\[
 \sigma(U)\in\{\eps,1,13\}.
\]
極値は連分数の交互辞書式順序で求める。この言語では、許された数字の後に2を
続ければ無限に延長できるので、貪欲な選択が正しい。
状態と偶奇が有限だから極値尾は最終的に周期的になる。$s=\sqrt{10}$ とすると尾の凸包は
\[
\begin{array}{c|cc}
\sigma&\min&\max\\\hline
\eps&-1+2s/5&-4/3+2s/3\\
1&-1/2+s/4&-4/3+2s/3\\
13&-1+2s/5&-5/3+2s/3
\end{array}
\]
である。一般の禁止語では無限延長できない状態を先に除く必要がある。
\texttt{generalized\_t.py} はこの処理を含む。

\section{全凸包を再帰させる方針への一様な障害}
\begin{proposition}\label{prop:obstruction}
$U,V$ がともに13で終わり、長さの偶奇が等しいとする。
$w_U=|\conv\KF(U3)|, w_V=|\conv\KF(V3)|$ とおく。
\[
 \frac56\le\frac{w_V}{w_U}\le\frac65
\]
なら、$H(U,V)$ の内部に $\KF(U)+\KF(V)$ と交わらない正長の開区間がある。
したがって、この条件を満たす全凸包を保持する無限の後続被覆は存在しない。
\end{proposition}
\begin{proof}
状態13の次の数字は2または3である。尾の3枝と2枝の凸包は順に $[a,b],[c,d]$ で、
\[
 a=-1+\frac{2s}{5},\quad b=\frac{5-s}{6},\quad
 c=\frac{s-1}{6},\quad d=\frac{2s-5}{3},\qquad a<b<c<d.
\]
接頭語の分母比を $r=q_{h-1}/q_h\in[0,1]$ とする。
$F_U$ で写した二枝の間の隙間の長さを $g_U$ とすると、偶奇にかかわらず
\[
 \frac{g_U}{w_U}=\frac{c-b}{b-a}\frac{1+ra}{1+rc}
 \ge\frac{c-b}{b-a}\frac{1+a}{1+c}
 =\frac{16}{9}-\frac{8s}{45}>\frac65.
\]
最後の不等式は $s<13/4$ から従う。$V$ にも同じ評価が成立する。
両長が偶数なら、3枝は下側にある。3枝同士の和の上端から
2枝・3枝の混合和の下端までの距離はそれぞれ
\[
 g_U-w_V>0,\qquad g_V-w_U>0.
\]
2枝同士の和はさらに上にあるから、両距離の小さい方が正長の隙間を与える。
両長が奇数なら順序が両側で反転するだけで、上側の3枝同士の和の直下に同じ隙間を得る。
各枝の端点は実際の $\KF$ の点なので、隙間は親の凸包の内部にある。
\end{proof}

最短の具体例は $U=V=13$ である。
\[
 H(13,13)=\left[\frac{44+4\sqrt{10}}{37},\frac{-8+4\sqrt{10}}3\right]
\]
に対し
\[
 [1.5358,1.5360]\subset H(13,13),\qquad
 [1.5358,1.5360]\cap(\KF(13)+\KF(13))=\varnothing.
\]
四つの子の上端・下端を $\mathbb Q(\sqrt{10})$ で比較するだけで検証できる。
これは\textbf{全体の $\KF+\KF$ の隙間だという主張ではない}。
許容族から該当する全凸包を除くか、追加禁止条件や担当区間を変える必要がある。

\section{有限探索の設計}
$F_U(x)=(p_{h-1}x+p_h)/(q_{h-1}x+q_h)$ より
\[
 F_U(x)-F_U(y)=\frac{(-1)^h(x-y)}{q_h^2(1+rx)(1+ry)}.
\]
形状 $r$ は末尾 $m$ 桁の逆向き連分数に $[1/4,4/5]$ を代入して囲う。
この前史区間は1,2,3の追加で不変である。実際の根は別途、形状の箱に入ることを確認する。
状態、偶奇、末尾 $m$ 桁を左右に保持する。
固定3桁では、末尾111の箱は $[14/23,9/14]$ で幅が $11/322$ 残る。
比だけを細分してもこの幅は減らない。
\texttt{--shape-tolerance} は、箱の幅が指定値を超える末尾だけを前方に一桁ずつ延長する。
例えば許容幅 $1/100$ では30種類の末尾（長さ2〜5）を使う。
偶奇を含めた60状態のすべてについて、子の形状の箱への包含を有理数で確認する。

アンカー $\zeta_U$ を尾の下端、または一定の周期点に選び、
\[
 D_U=|F'_U(\zeta_U)|,\qquad S=D_V/D_U
\]
を比の座標とする。$S$ を $[b^j,b^{j+1}]$ で分割する。試す補助区間は二方式である。
\begin{enumerate}
\item $H(U,V)=[l,h]$ の部分区間 $[l+p(h-l),l+q(h-l)]$。
\item 一定の $\zeta=(\sqrt5-1)/2$ に対し
\[
 [F_U(\zeta)+F_V(\zeta)+p(D_U+D_V),\quad
  F_U(\zeta)+F_V(\zeta)+q(D_U+D_V)].
\]
\end{enumerate}
第二方式は、禁止語で定義した一般化 T 区間そのものではない。
しかし、$p,q$ が有界、$S$ が正の上下界を持ち、合法な proper successor で閉じていれば、
定理\ref{thm:cover}と同じ極限論法が使える。両側が伸びることで $D_U+D_V\to0$ となり、
補助区間の中心と実際の極限和との差も0になるためである。
したがって、この方式では補助区間を各段階の凸包内に置く条件を省ける。

黄金比アンカーでは $F_1(\zeta)=\zeta$ だから、両側へ1を追加すると
\[
 D_{U1}=\zeta^2D_U,\qquad D_{V1}=\zeta^2D_V,\qquad S'=S.
\]
この等式は区間演算で曖昧にせず、代数的な一致として記録する。

各候補行について次を行う。
\begin{enumerate}
\item 許された追加語対を列挙し、子の形状と比の像を外側から囲う。
\item 像が触れる\emph{すべて}の比の箱で許容になっている子の区間型を採用する。
\item 二つの区間 $[l_i,h_i],[l_j,h_j]$ が交わることを
 $h_i\ge l_j$ と $h_j\ge l_i$ の\emph{両方}で箱全体にわたり確認する。
\item 親の両端を含む連結成分がない行を削除し、変化がなくなるまで反復する。
\end{enumerate}
比較は分母 $2^{48}$ の有理区間演算を用いる。入力の代数的数値は整数平方根によって
有理数で囲う。浮動小数点は候補順序や探索枝の削減だけに用い、採用の根拠には使わない。
探索では連結成分の保持候補数を制限しているため、空になったことから数学的な不可能性は結論しない。

\section{検証範囲と残る課題}
全接頭語版を直接扱う \texttt{search\_t.py} でも局所被覆を求めた。
$B_1=\{4,131,3\}$、$B_2=\{4,131\}$ を固定し、$I(U,V)=I(U,V;B_1,B_2)$ と略記すると
\[
 I(112,122)\subseteq I(1121,122)\cup I(1122,122)\cup I(1123,1222)
\]
が厳密に成立する。極値尾の周期表示を正規化し、同じ無限語は同一端点として扱い、
それ以外の比較は有理数による挟み撃ちで確認する。
ところが
\[
 H_*=[1.295458,1.295459]\subset I(1122,122),\qquad
 H_*\cap(\KF(1122)+\KF(122))=\varnothing.
\]
他の二つの子の補助区間も $H_*$ と交わらない。
さらに、根の左右の次の数字を一桁ずつ分けると、$H_*$ に触れる矩形の左追加数字は
すべて2である。従って、より強く
\[
 H_*\cap(\KF(112)+\KF(122))=\varnothing
\]
となる。この根の T 区間全体を出発区間にすること自体ができない。
局所的に三後続が見つかることと、無限に継続できることとの差を示す例である。
この結果は \texttt{verify\_local.py} が再検証する。
共有された議論の小区間 $[1.29288,1.292906]$ は $H_*$ と離れており、
本反例はその小区間の内点構成を否定しない。

禁止語の変更を有限範囲で全数検査すると、$B_0=\{4,131\}$、
$B_1=B_0\cup\{w\},B_2=B_0$、$1\le |w|\le3$ の39通りのうち、
3通りは空、残る36通りはすべて $H_*$ を含む。
両側に一語ずつ追加する場合も、重複を除いた780組のうち120組は空、657組は $H_*$ を含む。
残るのは追加語の組 $(3,3),(3,23),(23,23)$ であり、後二者も別の根の隙間を含む。
$(3,3)$、つまり両側二進型はこの検査では排除されなかったが、その包含は未証明である。
\texttt{verify\_forbidden\_family.py} がこれらを再検証する。

二進の極値そのものを保持する一様探索も行った。
F型、交換和T型、左右の二進型、両側二進型の五種類を使い、
各側に過去の3の出現フラグを持たせる。極値は $\mathbb Q(\sqrt3,\sqrt{10})$ で計算する。
末尾2桁の206,180行と、末尾3桁・比の公比 $51/50$ の6,496,720行はともに空になった。
部分区間型の最大の探索も、194,168,832行から20回の削除反復で空になった。
\textbf{131禁止については、内点の証明に必要な非空の閉じた許容族は得られていない。}

\texttt{verify\_obstructions.py} は命題\ref{prop:obstruction}と数値を指定した隙間を厳密に検証する。
別の禁止語オートマトンと有理挟み撃ちにより、長さ4以下の合法語について228個の
極値を照合した。空の言語、無限に延長できない状態、片方だけが空の交換項も確認した。
\texttt{verify\_t\_endpoints.py} はさらに1896個の極値と820個の空言語判定を別実装と照合する。
\texttt{check\_controls.py} は三つの検証器について、被覆できる人工例を受理し、
隙間・端点欠落を持つ人工例を拒否する。過去の違反フラグを含めた10件を確認した。
これらはプログラムの誤判定への検査であり、内点の証明を代替しない。

$\KF+\KF$ の内点の証明を完了するには、空でない閉じた許容表を得て、具体的な根の正長区間を
その表へ接続し、表を固定した再検証を通す必要がある。
有限段階で多数の行が残っていても、それだけでは証明にならない。
本稿の131禁止の探索は \texttt{stashed} の過去の証明書を読み込まず、Freiman の区間比較アルゴリズムを
131禁止に新しく適用したものである。
\end{document}
